\chapter*{References}
\addcontentsline{toc}{chapter}{References}

{[}1{]} O. Vinyals, A. Toshev, S. Bengio and D. Erhan, ``Show and tell: A neural image caption generator,'' in \emph{Proc. IEEE CVPR}, 2015, pp.~3156--3164.

{[}2{]} K. Xu, J. Ba, R. Kiros, K. Cho, A. Courville, R. Salakhutdinov, R. Zemel and Y. Bengio, ``Show, attend and tell: Neural image caption generation with visual attention,'' in \emph{Proc. ICML}, 2015, pp.~2048--2057.

{[}3{]} A. Vaswani, N. Shazeer, N. Parmar, J. Uszkoreit, L. Jones, A. N. Gomez, Ł. Kaiser and I. Polosukhin, ``Attention is all you need,'' in \emph{Advances in NeurIPS}, vol.~30, 2017.

{[}4{]} A. Dosovitskiy et al., ``An image is worth 16x16 words: Transformers for image recognition at scale,'' in \emph{Proc. ICLR}, 2021.

{[}5{]} Z. Liu, Y. Lin, Y. Cao, H. Hu, Y. Wei, Z. Zhang, S. Lin and B. Guo, ``Swin Transformer: Hierarchical vision transformer using shifted windows,'' in \emph{Proc. IEEE/CVF ICCV}, 2021, pp.~10012--10022.

{[}6{]} O. Pelka, S. Koitka, J. Rückert, F. Nensa and C. M. Friedrich, ``Radiology Objects in COntext (ROCO): A multimodal image dataset,'' in \emph{Intravascular Imaging and Computer Assisted Stenting and Large-Scale Annotation of Biomedical Data and Expert Label Synthesis}, LNCS vol.~11043, Springer, 2018, pp.~180--189.

{[}7{]} J. Rückert et al., ``ROCOv2: Radiology Objects in COntext version 2, an updated multimodal image dataset,'' \emph{Scientific Data}, vol.~11, art. 688, 2024.

{[}8{]} H. Lee, H. Cho, J. Park, J. Chae and J. Kim, ``Cross encoder-decoder transformer with global-local visual extractor for medical image captioning,'' \emph{Sensors}, vol.~22, no. 4, art. 1429, 2022.

{[}9{]} M. Pranay Kumar, V. Snigdha, R. Nandini and B. Indira Reddy, ``Image captioning generator using CNN and LSTM,'' \emph{International Journal for Research in Applied Science and Engineering Technology (IJRASET)}, vol.~10, no. 6, 2022, doi: 10.22214/ijraset.2022.44502.

{[}10{]} J. Deng, W. Dong, R. Socher, L.-J. Li, K. Li and L. Fei-Fei, ``ImageNet: A large-scale hierarchical image database,'' in \emph{Proc. IEEE CVPR}, 2009, pp.~248--255.

{[}11{]} K. He, X. Zhang, S. Ren and J. Sun, ``Deep residual learning for image recognition,'' in \emph{Proc. IEEE CVPR}, 2016, pp.~770--778.

{[}12{]} S. Hochreiter and J. Schmidhuber, ``Long short-term memory,'' \emph{Neural Computation}, vol.~9, no. 8, pp.~1735--1780, 1997.

{[}13{]} J. L. Ba, J. R. Kiros and G. E. Hinton, ``Layer normalization,'' arXiv:1607.06450, 2016.

{[}14{]} D. Hendrycks and K. Gimpel, ``Gaussian error linear units (GELUs),'' arXiv:1606.08415, 2016.

{[}15{]} G. Huang, Y. Sun, Z. Liu, D. Sedra and K. Q. Weinberger, ``Deep networks with stochastic depth,'' in \emph{Proc. ECCV}, 2016, pp.~646--661.

{[}16{]} O. Press and L. Wolf, ``Using the output embedding to improve language models,'' in \emph{Proc. EACL}, 2017, pp.~157--163.

{[}17{]} C. Szegedy, V. Vanhoucke, S. Ioffe, J. Shlens and Z. Wojna, ``Rethinking the Inception architecture for computer vision,'' in \emph{Proc. IEEE CVPR}, 2016, pp.~2818--2826.

{[}18{]} I. Loshchilov and F. Hutter, ``Decoupled weight decay regularization,'' in \emph{Proc. ICLR}, 2019.

{[}19{]} I. Loshchilov and F. Hutter, ``SGDR: Stochastic gradient descent with warm restarts,'' in \emph{Proc. ICLR}, 2017.

{[}20{]} Y. Wu et al., ``Google's neural machine translation system: Bridging the gap between human and machine translation,'' arXiv:1609.08144, 2016.

{[}21{]} K. Papineni, S. Roukos, T. Ward and W.-J. Zhu, ``BLEU: A method for automatic evaluation of machine translation,'' in \emph{Proc. ACL}, 2002, pp.~311--318.

{[}22{]} C.-Y. Lin, ``ROUGE: A package for automatic evaluation of summaries,'' in \emph{Text Summarization Branches Out (ACL Workshop)}, 2004, pp.~74--81.

{[}23{]} R. Vedantam, C. L. Zitnick and D. Parikh, ``CIDEr: Consensus-based image description evaluation,'' in \emph{Proc. IEEE CVPR}, 2015, pp.~4566--4575.

{[}24{]} S. Banerjee and A. Lavie, ``METEOR: An automatic metric for MT evaluation with improved correlation with human judgments,'' in \emph{Proc. ACL Workshop on Intrinsic and Extrinsic Evaluation Measures for MT and/or Summarization}, 2005, pp.~65--72.

{[}25{]} A. Paszke et al., ``PyTorch: An imperative style, high-performance deep learning library,'' in \emph{Advances in NeurIPS}, vol.~32, 2019.

{[}26{]} J. Rückert et al., ``Overview of ImageCLEFmedical 2023 -- Caption prediction and concept detection,'' in \emph{CLEF 2023 Working Notes}, CEUR Workshop Proceedings, vol.~3497, 2023.

{[}27{]} X. Mei et al., ``RadImageNet: An open radiologic deep learning research dataset for effective transfer learning,'' \emph{Radiology: Artificial Intelligence}, vol.~4, no. 5, e210315, 2022.

{[}28{]} Z. Wang, Z. Wu, D. Agarwal and J. Sun, ``MedCLIP: Contrastive learning from unpaired medical images and text,'' in \emph{Proc. EMNLP}, 2022, pp.~3876--3887.
